\documentclass[10pt,letterpaper]{article}
% assessment-curricula-common.tex
% Shared preamble for the software-security-assessment-modern-curricula suite.
% EXTRACTION PROVENANCE: the seven uploaded curriculum documents carried
% byte-identical inline preambles except for their \hypersetup{pdftitle=...}
% and \title{...} lines, which stay document-local. The lines below are the
% common lines (originally preamble lines 2-10 and 12-27) copied verbatim,
% in their original order. Documents load this file with
%   % assessment-curricula-common.tex
% Shared preamble for the software-security-assessment-modern-curricula suite.
% EXTRACTION PROVENANCE: the seven uploaded curriculum documents carried
% byte-identical inline preambles except for their \hypersetup{pdftitle=...}
% and \title{...} lines, which stay document-local. The lines below are the
% common lines (originally preamble lines 2-10 and 12-27) copied verbatim,
% in their original order. Documents load this file with
%   % assessment-curricula-common.tex
% Shared preamble for the software-security-assessment-modern-curricula suite.
% EXTRACTION PROVENANCE: the seven uploaded curriculum documents carried
% byte-identical inline preambles except for their \hypersetup{pdftitle=...}
% and \title{...} lines, which stay document-local. The lines below are the
% common lines (originally preamble lines 2-10 and 12-27) copied verbatim,
% in their original order. Documents load this file with
%   \input{assessment-curricula-common.tex}
% immediately after \documentclass, then set their own \hypersetup, \title,
% \author and \date.
\usepackage[T1]{fontenc}
\usepackage[utf8]{inputenc}
\usepackage{lmodern}
\usepackage[letterpaper,margin=0.80in,headheight=18pt,headsep=15pt]{geometry}
\usepackage{microtype,booktabs,longtable,tabularx,array,enumitem,xcolor,listings,tcolorbox,fancyhdr,titlesec,xurl,hyperref}
\definecolor{navy}{HTML}{18324E}
\definecolor{teal}{HTML}{136E77}
\definecolor{pale}{HTML}{EDF5F6}
\definecolor{papergray}{HTML}{F4F6F8}
\titleformat{\section}{\Large\bfseries\color{navy}}{\thesection}{.65em}{}
\titleformat{\subsection}{\large\bfseries\color{teal}}{\thesubsection}{.65em}{}
\setcounter{tocdepth}{2}
\setlength{\parskip}{4pt}
\setlength{\parindent}{0pt}
\emergencystretch=2em
\pagestyle{fancy}
\fancyhf{}
\fancyhead[L]{\footnotesize\color{navy}Software Security Assessment | Curriculum Suite}
\fancyhead[R]{\footnotesize\color{teal}2026 edition}
\fancyfoot[L]{\footnotesize Educational use | Authorized local laboratories}
\fancyfoot[R]{\footnotesize\thepage}
\lstdefinestyle{reviewcode}{basicstyle=\ttfamily\footnotesize,breaklines=true,columns=fullflexible,keepspaces=true,showstringspaces=false,frame=single,rulecolor=\color{teal},backgroundcolor=\color{papergray},aboveskip=8pt,belowskip=8pt}
\lstset{style=reviewcode}
\newtcolorbox{insight}[1][]{colback=pale,colframe=teal,boxrule=.5pt,arc=2mm,#1}
\newtcolorbox{historynote}[1][]{colback=papergray,colframe=navy,boxrule=.5pt,arc=2mm,#1}

% immediately after \documentclass, then set their own \hypersetup, \title,
% \author and \date.
\usepackage[T1]{fontenc}
\usepackage[utf8]{inputenc}
\usepackage{lmodern}
\usepackage[letterpaper,margin=0.80in,headheight=18pt,headsep=15pt]{geometry}
\usepackage{microtype,booktabs,longtable,tabularx,array,enumitem,xcolor,listings,tcolorbox,fancyhdr,titlesec,xurl,hyperref}
\definecolor{navy}{HTML}{18324E}
\definecolor{teal}{HTML}{136E77}
\definecolor{pale}{HTML}{EDF5F6}
\definecolor{papergray}{HTML}{F4F6F8}
\titleformat{\section}{\Large\bfseries\color{navy}}{\thesection}{.65em}{}
\titleformat{\subsection}{\large\bfseries\color{teal}}{\thesubsection}{.65em}{}
\setcounter{tocdepth}{2}
\setlength{\parskip}{4pt}
\setlength{\parindent}{0pt}
\emergencystretch=2em
\pagestyle{fancy}
\fancyhf{}
\fancyhead[L]{\footnotesize\color{navy}Software Security Assessment | Curriculum Suite}
\fancyhead[R]{\footnotesize\color{teal}2026 edition}
\fancyfoot[L]{\footnotesize Educational use | Authorized local laboratories}
\fancyfoot[R]{\footnotesize\thepage}
\lstdefinestyle{reviewcode}{basicstyle=\ttfamily\footnotesize,breaklines=true,columns=fullflexible,keepspaces=true,showstringspaces=false,frame=single,rulecolor=\color{teal},backgroundcolor=\color{papergray},aboveskip=8pt,belowskip=8pt}
\lstset{style=reviewcode}
\newtcolorbox{insight}[1][]{colback=pale,colframe=teal,boxrule=.5pt,arc=2mm,#1}
\newtcolorbox{historynote}[1][]{colback=papergray,colframe=navy,boxrule=.5pt,arc=2mm,#1}

% immediately after \documentclass, then set their own \hypersetup, \title,
% \author and \date.
\usepackage[T1]{fontenc}
\usepackage[utf8]{inputenc}
\usepackage{lmodern}
\usepackage[letterpaper,margin=0.80in,headheight=18pt,headsep=15pt]{geometry}
\usepackage{microtype,booktabs,longtable,tabularx,array,enumitem,xcolor,listings,tcolorbox,fancyhdr,titlesec,xurl,hyperref}
\definecolor{navy}{HTML}{18324E}
\definecolor{teal}{HTML}{136E77}
\definecolor{pale}{HTML}{EDF5F6}
\definecolor{papergray}{HTML}{F4F6F8}
\titleformat{\section}{\Large\bfseries\color{navy}}{\thesection}{.65em}{}
\titleformat{\subsection}{\large\bfseries\color{teal}}{\thesubsection}{.65em}{}
\setcounter{tocdepth}{2}
\setlength{\parskip}{4pt}
\setlength{\parindent}{0pt}
\emergencystretch=2em
\pagestyle{fancy}
\fancyhf{}
\fancyhead[L]{\footnotesize\color{navy}Software Security Assessment | Curriculum Suite}
\fancyhead[R]{\footnotesize\color{teal}2026 edition}
\fancyfoot[L]{\footnotesize Educational use | Authorized local laboratories}
\fancyfoot[R]{\footnotesize\thepage}
\lstdefinestyle{reviewcode}{basicstyle=\ttfamily\footnotesize,breaklines=true,columns=fullflexible,keepspaces=true,showstringspaces=false,frame=single,rulecolor=\color{teal},backgroundcolor=\color{papergray},aboveskip=8pt,belowskip=8pt}
\lstset{style=reviewcode}
\newtcolorbox{insight}[1][]{colback=pale,colframe=teal,boxrule=.5pt,arc=2mm,#1}
\newtcolorbox{historynote}[1][]{colback=papergray,colframe=navy,boxrule=.5pt,arc=2mm,#1}

\hypersetup{colorlinks=true,linkcolor=navy,urlcolor=teal,pdftitle={Windows Software Security Assessment}}
\title{\vspace{-1.6cm}\Huge\color{navy}\bfseries Windows Software Security Assessment\\[8pt]\large\color{teal}A 16-week standalone, evidence-led professional learning pathway}
\author{Software Security Engineering Curriculum | Independent instructional design}
\date{Version 1.0 | 17 October 2026}
\begin{document}
\maketitle
\begin{center}\small Based on the learning tracks in Dowd, McDonald and Schuh (2007);\
contemporary topics are separately labeled as instructional extensions.\end{center}
\begin{insight}\textbf{Audience and ethics.} This curriculum is for authorized software assessment only. All exercises use instructor-owned source, synthetic input, offline captures, or isolated localhost/VM environments. No third-party probing is required.\end{insight}
\tableofcontents
\clearpage
\section{Executive overview and educational scope}
This independent Windows curriculum prepares experienced developers and emerging security reviewers to investigate software defects as potential breaches of explicit security expectations. It follows the original book's common Chapter 1--4 foundation followed by Chapters 5, 6, 7, 8, 11, 12, 13 in their preface order. Work prioritizes manual auditing and verification before automated classification.

Ideal audience: application security engineers, secure developers, auditors and graduate cybersecurity learners. Expected entry skills: read one programming language confidently, understand variables/functions/HTTP or system APIs where relevant, use Git, and write elementary tests. Core laboratory languages: C, C++, PowerShell, C\#. Specialized local execution environment: isolated Windows evaluation VM; Visual Studio Build Tools, WinDbg, PowerShell.

\subsection{Professional learning outcomes}
\begin{itemize}[leftmargin=*,itemsep=2pt]
\item Evaluate Windows application architecture, privileges and untrusted data paths against explicit security invariants.
\item Conduct independent structured source review, characterize preconditions and avoid unsupported exploitability claims.
\item Validate a suspect flaw with reproducible safe tests, fix the root cause, and demonstrate regression coverage.
\item Reconcile modern scanner evidence with manual reachability reasoning and report uncertainty accurately.
\item Deliver a prioritized audit report and defended residual-risk decision in an engineering change process.
\end{itemize}
\section{Original textbook alignment and source provenance}
The source-derived sequence is the preface selection, not an invented modern order. The book recommends Chapters 1--4 as common prerequisites. Printed textbook page references below are distinct from PDF viewer page indices. Supplemental examples labeled as contemporary extensions do not appear in the 2007 source.

\begingroup\footnotesize\setlength\tabcolsep{3pt}
\begin{longtable}{>{\raggedright\arraybackslash}p{\dimexpr0.310\textwidth-6pt\relax}>{\raggedright\arraybackslash}p{\dimexpr0.370\textwidth-6pt\relax}>{\raggedright\arraybackslash}p{\dimexpr0.320\textwidth-6pt\relax}}\toprule
\textbf{Chapter and printed pages} & \textbf{Textbook-grounded subject} & \textbf{Modern instructional bridge}\\\midrule\endfirsthead
\toprule \textbf{Chapter and printed pages} & \textbf{Textbook-grounded subject} & \textbf{Modern instructional bridge}\\\midrule\endhead
1: Software Vulnerability Fundamentals (pp. 3--24) & vulnerability classification; security policies; trust relationships; audit vs. black-box testing & CWE and risk-informed prioritization\\\midrule
2: Design Review (pp. 25--66) & architecture, trust modeling, authentication, authorization and threat identification & STRIDE and cloud-native trust boundaries\\\midrule
3: Operational Review (pp. 67--90) & attack-surface assessment, insecure defaults, secure channels and protective measures & containers, IAM and configuration-as-code\\\midrule
4: Application Review Process (pp. 91--166) & scoping, code navigation, candidate points, data flow, desk-checking and reporting & CodeQL, SARIF, SAST triage and CI verification\\\midrule
5: Memory Corruption (pp. 167--202) & stack/heap overflows, off-by-one errors, mitigations and impact assessment & sanitizers, fuzzing and memory-safe implementation alternatives\\\midrule
6: C Language Issues (pp. 203--296) & integer boundaries, type conversion, promotions, pointers and undefined assumptions & checked arithmetic and contemporary C/C++ toolchains\\\midrule
7: Program Building Blocks (pp. 297--386) & variable state, control flow, return values, memory allocation and double frees & semantic data flow and ownership/static analysis\\\midrule
8: Strings and Metacharacters (pp. 387--458) & strings, NUL, encodings, path/shell metacharacters and injection & contextual encoding, parameterization and Unicode normalization\\\midrule
11: Windows I: Objects and the File System (pp. 625--684) & Windows handles, access tokens, security descriptors, filesystem and registry & modern Windows service identities and DLL search safety\\\midrule
12: Windows II: Interprocess Communication (pp. 685--754) & pipes, impersonation, window messaging, RPC and COM/DCOM & named-pipe access checks and Windows RPC hardening\\\midrule
13: Synchronization and State (pp. 755--828) & race conditions, signals, reentrancy, deadlocks and threads & thread sanitizers, asynchronous services and atomicity\\\midrule
\bottomrule\end{longtable}\endgroup
\begin{historynote}\textbf{Historical interpretation.} Legacy CGI, ActiveX, obsolete APIs, pre-modern mitigation assumptions and earlier HTTP frameworks are studied for their vulnerability mechanisms. They must not be represented as current default production guidance. Modern platform behavior is checked against vendor specifications.\end{historynote}
\section{Sixteen-week course schedule}
Workload: 8 hours/week (2 guided instruction, 2 directed reading/review, 4 practical lab or assessment). Weeks 1--4 deliver the shared foundation. Weeks 5--14 follow the specialization sequence without reordering the book chapters. Weeks 15--16 cover full capstone synthesis.

\begingroup\footnotesize\setlength\tabcolsep{3pt}
\begin{longtable}{>{\raggedright\arraybackslash}p{\dimexpr0.065\textwidth-6pt\relax}>{\raggedright\arraybackslash}p{\dimexpr0.390\textwidth-6pt\relax}>{\raggedright\arraybackslash}p{\dimexpr0.120\textwidth-6pt\relax}>{\raggedright\arraybackslash}p{\dimexpr0.425\textwidth-6pt\relax}}\toprule
\textbf{Week} & \textbf{Module and learning focus} & \textbf{Source} & \textbf{Primary evidence}\\\midrule\endfirsthead
\toprule \textbf{Week} & \textbf{Module and learning focus} & \textbf{Source} & \textbf{Primary evidence}\\\midrule\endhead
1 & F1: Security policies and vulnerability fundamentals & Ch. 1 & classification.csv; 1-page policy assumptions; issue evidence notes\\\midrule
2 & F2: Architecture, trust and threat modeling & Ch. 2 & DFD; threat register; ranked audit entry points\\\midrule
3 & F3: Operational exposure and deployment review & Ch. 3 & exposure matrix; hardened configuration diff; verification log\\\midrule
4 & F4: Manual code auditing and evidence standards & Ch. 4 & annotated flow trace; triage record; remediation PR with regression tests\\\midrule
5 & W1: Native buffer and allocation safety & Ch. 5 & memory map; corrected code; validation cases\\\midrule
6 & W2: Conversions and native API lengths & Ch. 6 & conversion diagram; regression suite; code diff\\\midrule
7 & W3: Ownership, handles and failure returns & Ch. 7 & path coverage table; RAII rewrite; static report\\\midrule
8 & W4: Unicode, filenames and dynamic execution & Ch. 8 & string transformation map; patched helper; tests\\\midrule
9 & W5: Tokens and object access controls & Ch. 11 & ACL matrix; token reasoning; PowerShell verification output\\\midrule
10 & W6: Files, registry and DLL loading & Ch. 11 & load trace; ACL before/after; mitigation evidence\\\midrule
11 & W7: Named pipes and impersonation & Ch. 12 & pipe trust diagram; fixed checks; access test\\\midrule
12 & W8: RPC, COM and service boundaries & Ch. 12 & call graph; authorization matrix; test protocol\\\midrule
13 & W9: Windows thread safety and races & Ch. 13 & interleaving chart; lock protocol; tests\\\midrule
14 & W10: Windows native audit integration & Ch. 13 & consolidated report; test runs; residual-risk memo\\\midrule
15 & C1: Capstone: audit and evidence & Ch. 4 & scope; architecture; threat model; annotated findings; controlled validation logs\\\midrule
16 & C2: Capstone: remediation and assurance & Ch. 4 & patchset; test evidence; SARIF/triage record; final written report and presentation\\\midrule
\bottomrule\end{longtable}\endgroup
\subsection{Eight- and 24-week delivery variants}
\begin{itemize}[leftmargin=*,itemsep=2pt]
\item Accelerated 8-week cohort: teach foundation F1+F2 in week 1 and F3+F4 in week 2; combine specialization labs by adjacent chapter in weeks 3--6; preserve each original chapter outcome; reserve weeks 7--8 for capstone. Plan 16 hours per week, not reduced coverage.
\item Extended 24-week cohort: use one week per foundation module; spread specialization labs over weeks 5--20 with extra code review and independent reading; capstone and peer review occupy weeks 21--24. Plan 6--8 hours/week.
\item For compressed delivery, explicitly list dependent skill checkpoints and require remediation of failed prerequisites before learners proceed.
\end{itemize}
\section{Detailed instructional modules and executable assessment protocols}
The code miniatures in this section are original, compact review exercises. They illustrate an isolated concept, not a full buildable application by themselves; the lab procedure states the platform and evidence required for the local application. Instructors provide pinned lab targets or use the companion starter fixtures. Security verification must be performed before marking outcomes complete.

\subsection{Week 01 | F1 | Security policies and vulnerability fundamentals}
\textbf{Prerequisite:} programming fundamentals\hfill\textbf{Time:} 2h instruction + 2h reading + 4h lab

\textbf{Source reading [textbook-grounded]:} Ch. 1, Software Vulnerability Fundamentals, pp. 3--24; focus: Vulnerabilities; Security Policies; Classifying Vulnerabilities.

\textbf{Contemporary extension [not from textbook]:} CWE, CVSS 4.0, EPSS and Secure by Design

\textbf{Competency objectives (analyze/evaluate/create):}
\begin{itemize}[leftmargin=*,itemsep=2pt]
\item Diagram the input, privileged decision or state transition implicated in security policies and vulnerability fundamentals.
\item Audit the nominated source operations and explain security conditions for create an evidence-based classification for six local code excerpts.
\item Design a correction with both a positive test and an adversarial negative case; defend the residual-risk conclusion.
\end{itemize}
\textbf{Lecture and instructor demonstration.} Teach the relevant mechanism from Chapter 1, including the interaction between input representation, trust relationships, control flow and impacts. Instructor narrates a source-level trace, identifies a missing invariant, predicts a failing test, and demonstrates a corrected control. Link the older mechanism to CWE, CVSS 4.0, EPSS and Secure by Design, making version differences explicit.

\begin{insight}\textbf{Lab F1:} Create an evidence-based classification for six local code excerpts.\\\textbf{Environment:} isolated Windows evaluation VM; Visual Studio Build Tools, WinDbg, PowerShell\\\textbf{Security assumption:} The sample is local/owned and untrusted inputs cannot be assumed canonical, authorized, or correctly sized.\end{insight}
\textbf{Starter code review miniature (new instructional example; not textbook text):}
\begin{lstlisting}[language=Python]
# Review question: which security expectation is violated?
def read_invoice(db, caller, invoice_id):
    invoice = db.lookup(invoice_id)
    return invoice  # Missing ownership check
# Safer: require invoice.owner_id == caller.id before return
\end{lstlisting}
\textbf{Laboratory execution procedure}
\begin{itemize}[leftmargin=*,itemsep=2pt]
\item Prepare a fresh working copy. Save the miniature as a discussion fixture, or map its flaw to the provided local application. Record baseline behavior and tool versions.
\item Formulate the security invariant and draw a source-to-sink or state-transition trace. List exactly which inputs, principals or concurrent transitions are under adversarial control.
\item Run controlled verification: Classify design vs. implementation vs. operational error; articulate the violated security expectation.
\item Patch the underlying defect using a guard, authorization check, atomic operation, API contract or safe construction appropriate to the platform; note defense-in-depth controls separately.
\item Rerun the original case plus a valid control, an invalid case and a relevant boundary case. Save executable test input, actual output and review notes.
\item Export the weekly artifacts: classification.csv; 1-page policy assumptions; issue evidence notes.
\end{itemize}
\textbf{Knowledge check.} Which security assumption was false? What exact precondition turns the programming error into a security vulnerability? What evidence would refute the finding?

\textbf{Scoring (10 points).} 3 root-cause/trace, 2 adversary conditions, 2 correction, 2 executable before/after evidence, 1 limits and clarity.

\textbf{Exit criteria.} A reviewer can reconstruct the flaw and confirm the relevant invariant now holds using the supplied tests; uncertain severity and environmental dependencies are recorded.

\textbf{Common pitfall.} Incorrectly treating every bug as an exploitable vulnerability.

\textbf{Next dependency.} Transfer the acquired audit pattern into the next original textbook chapter or the capstone; do not advance if trace, fix or evidence requirements are missing.

\subsection{Week 02 | F2 | Architecture, trust and threat modeling}
\textbf{Prerequisite:} F1 completion\hfill\textbf{Time:} 2h instruction + 2h reading + 4h lab

\textbf{Source reading [textbook-grounded]:} Ch. 2, Design Review, pp. 25--66; focus: Trust Relationships; Threat Modeling; Architecture Modeling.

\textbf{Contemporary extension [not from textbook]:} STRIDE, data-flow diagrams, security invariants

\textbf{Competency objectives (analyze/evaluate/create):}
\begin{itemize}[leftmargin=*,itemsep=2pt]
\item Diagram the input, privileged decision or state transition implicated in architecture, trust and threat modeling.
\item Audit the nominated source operations and explain security conditions for model a local two-tier service with authentication and storage boundaries.
\item Design a correction with both a positive test and an adversarial negative case; defend the residual-risk conclusion.
\end{itemize}
\textbf{Lecture and instructor demonstration.} Teach the relevant mechanism from Chapter 2, including the interaction between input representation, trust relationships, control flow and impacts. Instructor narrates a source-level trace, identifies a missing invariant, predicts a failing test, and demonstrates a corrected control. Link the older mechanism to STRIDE, data-flow diagrams, security invariants, making version differences explicit.

\begin{insight}\textbf{Lab F2:} Model a local two-tier service with authentication and storage boundaries.\\\textbf{Environment:} isolated Windows evaluation VM; Visual Studio Build Tools, WinDbg, PowerShell\\\textbf{Security assumption:} The sample is local/owned and untrusted inputs cannot be assumed canonical, authorized, or correctly sized.\end{insight}
\textbf{Starter code review miniature (new instructional example; not textbook text):}
\begin{lstlisting}[language=Python]
# Request arrives from an untrusted client.
def get_report(subject, report):
    if not subject.authenticated:
        raise PermissionError("login required")
    return report  # Missing object-level authorization
# Enforce: subject.may_read(report) before returning
\end{lstlisting}
\textbf{Laboratory execution procedure}
\begin{itemize}[leftmargin=*,itemsep=2pt]
\item Prepare a fresh working copy. Save the miniature as a discussion fixture, or map its flaw to the provided local application. Record baseline behavior and tool versions.
\item Formulate the security invariant and draw a source-to-sink or state-transition trace. List exactly which inputs, principals or concurrent transitions are under adversarial control.
\item Run controlled verification: Identify assets, attackers, flows and an authorization bypass route in the diagram.
\item Patch the underlying defect using a guard, authorization check, atomic operation, API contract or safe construction appropriate to the platform; note defense-in-depth controls separately.
\item Rerun the original case plus a valid control, an invalid case and a relevant boundary case. Save executable test input, actual output and review notes.
\item Export the weekly artifacts: DFD; threat register; ranked audit entry points.
\end{itemize}
\textbf{Knowledge check.} Which security assumption was false? What exact precondition turns the programming error into a security vulnerability? What evidence would refute the finding?

\textbf{Scoring (10 points).} 3 root-cause/trace, 2 adversary conditions, 2 correction, 2 executable before/after evidence, 1 limits and clarity.

\textbf{Exit criteria.} A reviewer can reconstruct the flaw and confirm the relevant invariant now holds using the supplied tests; uncertain severity and environmental dependencies are recorded.

\textbf{Common pitfall.} Confusing authentication with permission to access a resource.

\textbf{Next dependency.} Transfer the acquired audit pattern into the next original textbook chapter or the capstone; do not advance if trace, fix or evidence requirements are missing.

\subsection{Week 03 | F3 | Operational exposure and deployment review}
\textbf{Prerequisite:} F2 completion\hfill\textbf{Time:} 2h instruction + 2h reading + 4h lab

\textbf{Source reading [textbook-grounded]:} Ch. 3, Operational Review, pp. 67--90; focus: Exposure; Insecure Defaults; Protective Measures.

\textbf{Contemporary extension [not from textbook]:} container settings, least-privilege IAM and secure defaults

\textbf{Competency objectives (analyze/evaluate/create):}
\begin{itemize}[leftmargin=*,itemsep=2pt]
\item Diagram the input, privileged decision or state transition implicated in operational exposure and deployment review.
\item Audit the nominated source operations and explain security conditions for inspect a prepared Compose manifest and reverse-proxy configuration.
\item Design a correction with both a positive test and an adversarial negative case; defend the residual-risk conclusion.
\end{itemize}
\textbf{Lecture and instructor demonstration.} Teach the relevant mechanism from Chapter 3, including the interaction between input representation, trust relationships, control flow and impacts. Instructor narrates a source-level trace, identifies a missing invariant, predicts a failing test, and demonstrates a corrected control. Link the older mechanism to container settings, least-privilege IAM and secure defaults, making version differences explicit.

\begin{insight}\textbf{Lab F3:} Inspect a prepared Compose manifest and reverse-proxy configuration.\\\textbf{Environment:} isolated Windows evaluation VM; Visual Studio Build Tools, WinDbg, PowerShell\\\textbf{Security assumption:} The sample is local/owned and untrusted inputs cannot be assumed canonical, authorized, or correctly sized.\end{insight}
\textbf{Starter code review miniature (new instructional example; not textbook text):}
\begin{lstlisting}[language=Python]
# Operational policy fixture; do not deploy this configuration
service = {"listen": "0.0.0.0:8080", "debug": True,
           "run_as": "root"}
# Review ingress, error disclosure, service identity and secrets.
\end{lstlisting}
\textbf{Laboratory execution procedure}
\begin{itemize}[leftmargin=*,itemsep=2pt]
\item Prepare a fresh working copy. Save the miniature as a discussion fixture, or map its flaw to the provided local application. Record baseline behavior and tool versions.
\item Formulate the security invariant and draw a source-to-sink or state-transition trace. List exactly which inputs, principals or concurrent transitions are under adversarial control.
\item Run controlled verification: Identify unnecessary ports, unsafe environment secrets and excessive process privileges.
\item Patch the underlying defect using a guard, authorization check, atomic operation, API contract or safe construction appropriate to the platform; note defense-in-depth controls separately.
\item Rerun the original case plus a valid control, an invalid case and a relevant boundary case. Save executable test input, actual output and review notes.
\item Export the weekly artifacts: exposure matrix; hardened configuration diff; verification log.
\end{itemize}
\textbf{Knowledge check.} Which security assumption was false? What exact precondition turns the programming error into a security vulnerability? What evidence would refute the finding?

\textbf{Scoring (10 points).} 3 root-cause/trace, 2 adversary conditions, 2 correction, 2 executable before/after evidence, 1 limits and clarity.

\textbf{Exit criteria.} A reviewer can reconstruct the flaw and confirm the relevant invariant now holds using the supplied tests; uncertain severity and environmental dependencies are recorded.

\textbf{Common pitfall.} Assuming a network control repairs an application logic defect.

\textbf{Next dependency.} Transfer the acquired audit pattern into the next original textbook chapter or the capstone; do not advance if trace, fix or evidence requirements are missing.

\subsection{Week 04 | F4 | Manual code auditing and evidence standards}
\textbf{Prerequisite:} F3 completion\hfill\textbf{Time:} 2h instruction + 2h reading + 4h lab

\textbf{Source reading [textbook-grounded]:} Ch. 4, Application Review Process, pp. 91--166; focus: Code-Auditing Strategies; Internal Flow Analysis; Reporting.

\textbf{Contemporary extension [not from textbook]:} CodeQL, Semgrep, SARIF 2.1.0 and CI check runs

\textbf{Competency objectives (analyze/evaluate/create):}
\begin{itemize}[leftmargin=*,itemsep=2pt]
\item Diagram the input, privileged decision or state transition implicated in manual code auditing and evidence standards.
\item Audit the nominated source operations and explain security conditions for trace a supplied input through normalization, permission checks and an unsafe sink.
\item Design a correction with both a positive test and an adversarial negative case; defend the residual-risk conclusion.
\end{itemize}
\textbf{Lecture and instructor demonstration.} Teach the relevant mechanism from Chapter 4, including the interaction between input representation, trust relationships, control flow and impacts. Instructor narrates a source-level trace, identifies a missing invariant, predicts a failing test, and demonstrates a corrected control. Link the older mechanism to CodeQL, Semgrep, SARIF 2.1.0 and CI check runs, making version differences explicit.

\begin{insight}\textbf{Lab F4:} Trace a supplied input through normalization, permission checks and an unsafe sink.\\\textbf{Environment:} isolated Windows evaluation VM; Visual Studio Build Tools, WinDbg, PowerShell\\\textbf{Security assumption:} The sample is local/owned and untrusted inputs cannot be assumed canonical, authorized, or correctly sized.\end{insight}
\textbf{Starter code review miniature (new instructional example; not textbook text):}
\begin{lstlisting}[language=Python]
def triage(source, sink, guard):
    # A finding requires more than matching a dangerous call.
    return {"source": source, "sink": sink,
            "guard": guard, "reachable": None,
            "validated": False}
# Auditor must investigate real guards and execution paths.
\end{lstlisting}
\textbf{Laboratory execution procedure}
\begin{itemize}[leftmargin=*,itemsep=2pt]
\item Prepare a fresh working copy. Save the miniature as a discussion fixture, or map its flaw to the provided local application. Record baseline behavior and tool versions.
\item Formulate the security invariant and draw a source-to-sink or state-transition trace. List exactly which inputs, principals or concurrent transitions are under adversarial control.
\item Run controlled verification: Distinguish reachability, exploitable conditions and scanner confidence before proposing a fix.
\item Patch the underlying defect using a guard, authorization check, atomic operation, API contract or safe construction appropriate to the platform; note defense-in-depth controls separately.
\item Rerun the original case plus a valid control, an invalid case and a relevant boundary case. Save executable test input, actual output and review notes.
\item Export the weekly artifacts: annotated flow trace; triage record; remediation PR with regression tests.
\end{itemize}
\textbf{Knowledge check.} Which security assumption was false? What exact precondition turns the programming error into a security vulnerability? What evidence would refute the finding?

\textbf{Scoring (10 points).} 3 root-cause/trace, 2 adversary conditions, 2 correction, 2 executable before/after evidence, 1 limits and clarity.

\textbf{Exit criteria.} A reviewer can reconstruct the flaw and confirm the relevant invariant now holds using the supplied tests; uncertain severity and environmental dependencies are recorded.

\textbf{Common pitfall.} Accepting a scanner alert without following the relevant execution path.

\textbf{Next dependency.} Transfer the acquired audit pattern into the next original textbook chapter or the capstone; do not advance if trace, fix or evidence requirements are missing.

\begin{historynote}\textbf{Checkpoint:} Foundation exit examination: annotate one untrusted-data flow and defend a finding.\end{historynote}
\subsection{Week 05 | W1 | Native buffer and allocation safety}
\textbf{Prerequisite:} F1--F4; preceding specialized audit lab\hfill\textbf{Time:} 2h instruction + 2h reading + 4h lab

\textbf{Source reading [textbook-grounded]:} Ch. 5, Memory Corruption, pp. 167--202; focus: Buffer Overflows; Protection Mechanisms.

\textbf{Contemporary extension [not from textbook]:} Windows ASan, CFG and modern compiler switches

\textbf{Competency objectives (analyze/evaluate/create):}
\begin{itemize}[leftmargin=*,itemsep=2pt]
\item Diagram the input, privileged decision or state transition implicated in native buffer and allocation safety.
\item Audit the nominated source operations and explain security conditions for assess unchecked array copy in a local C++ component.
\item Design a correction with both a positive test and an adversarial negative case; defend the residual-risk conclusion.
\end{itemize}
\textbf{Lecture and instructor demonstration.} Teach the relevant mechanism from Chapter 5, including the interaction between input representation, trust relationships, control flow and impacts. Instructor narrates a source-level trace, identifies a missing invariant, predicts a failing test, and demonstrates a corrected control. Link the older mechanism to Windows ASan, CFG and modern compiler switches, making version differences explicit.

\begin{insight}\textbf{Lab W1:} Assess unchecked array copy in a local C++ component.\\\textbf{Environment:} isolated Windows evaluation VM; Visual Studio Build Tools, WinDbg, PowerShell\\\textbf{Security assumption:} The sample is local/owned and untrusted inputs cannot be assumed canonical, authorized, or correctly sized.\end{insight}
\textbf{Starter code review miniature (new instructional example; not textbook text):}
\begin{lstlisting}[language=C]
/* Example is intentionally flawed: local demonstration only. */
void read_field(const char *src, size_t n) {
    char dst[32];
    if (n <= sizeof dst) memcpy(dst, src, n);
    /* If string semantics require NUL, capacity differs. */
}
/* Review actual source-buffer length too. */
\end{lstlisting}
\textbf{Laboratory execution procedure}
\begin{itemize}[leftmargin=*,itemsep=2pt]
\item Prepare a fresh working copy. Save the miniature as a discussion fixture, or map its flaw to the provided local application. Record baseline behavior and tool versions.
\item Formulate the security invariant and draw a source-to-sink or state-transition trace. List exactly which inputs, principals or concurrent transitions are under adversarial control.
\item Run controlled verification: Evaluate boundary inputs under available sanitizers; explain that CFG does not repair an overflow.
\item Patch the underlying defect using a guard, authorization check, atomic operation, API contract or safe construction appropriate to the platform; note defense-in-depth controls separately.
\item Rerun the original case plus a valid control, an invalid case and a relevant boundary case. Save executable test input, actual output and review notes.
\item Export the weekly artifacts: memory map; corrected code; validation cases.
\end{itemize}
\textbf{Knowledge check.} Which security assumption was false? What exact precondition turns the programming error into a security vulnerability? What evidence would refute the finding?

\textbf{Scoring (10 points).} 3 root-cause/trace, 2 adversary conditions, 2 correction, 2 executable before/after evidence, 1 limits and clarity.

\textbf{Exit criteria.} A reviewer can reconstruct the flaw and confirm the relevant invariant now holds using the supplied tests; uncertain severity and environmental dependencies are recorded.

\textbf{Common pitfall.} Equating exploit-mitigation deployment with vulnerability elimination.

\textbf{Next dependency.} Transfer the acquired audit pattern into the next original textbook chapter or the capstone; do not advance if trace, fix or evidence requirements are missing.

\subsection{Week 06 | W2 | Conversions and native API lengths}
\textbf{Prerequisite:} F1--F4; preceding specialized audit lab\hfill\textbf{Time:} 2h instruction + 2h reading + 4h lab

\textbf{Source reading [textbook-grounded]:} Ch. 6, C Language Issues, pp. 203--296; focus: Signed/Unsigned Conversions; Truncation.

\textbf{Contemporary extension [not from textbook]:} checked arithmetic and Windows integer types

\textbf{Competency objectives (analyze/evaluate/create):}
\begin{itemize}[leftmargin=*,itemsep=2pt]
\item Diagram the input, privileged decision or state transition implicated in conversions and native api lengths.
\item Audit the nominated source operations and explain security conditions for audit DWORD and size\_t conversions near a buffer allocation.
\item Design a correction with both a positive test and an adversarial negative case; defend the residual-risk conclusion.
\end{itemize}
\textbf{Lecture and instructor demonstration.} Teach the relevant mechanism from Chapter 6, including the interaction between input representation, trust relationships, control flow and impacts. Instructor narrates a source-level trace, identifies a missing invariant, predicts a failing test, and demonstrates a corrected control. Link the older mechanism to checked arithmetic and Windows integer types, making version differences explicit.

\begin{insight}\textbf{Lab W2:} Audit DWORD and size\_t conversions near a buffer allocation.\\\textbf{Environment:} isolated Windows evaluation VM; Visual Studio Build Tools, WinDbg, PowerShell\\\textbf{Security assumption:} The sample is local/owned and untrusted inputs cannot be assumed canonical, authorized, or correctly sized.\end{insight}
\textbf{Starter code review miniature (new instructional example; not textbook text):}
\begin{lstlisting}[language=C]
/* Risk: negative signed length becomes huge size_t. */
int declared = read_signed_length();
size_t n = (size_t)declared;
char *p = malloc(n);
/* Correct: check declared >= 0, upper bound, and math
   BEFORE cast/allocation; test allocation failure. */
\end{lstlisting}
\textbf{Laboratory execution procedure}
\begin{itemize}[leftmargin=*,itemsep=2pt]
\item Prepare a fresh working copy. Save the miniature as a discussion fixture, or map its flaw to the provided local application. Record baseline behavior and tool versions.
\item Formulate the security invariant and draw a source-to-sink or state-transition trace. List exactly which inputs, principals or concurrent transitions are under adversarial control.
\item Run controlled verification: Test boundary lengths and the behavior of a 32-bit-to-64-bit conversion with a portable harness.
\item Patch the underlying defect using a guard, authorization check, atomic operation, API contract or safe construction appropriate to the platform; note defense-in-depth controls separately.
\item Rerun the original case plus a valid control, an invalid case and a relevant boundary case. Save executable test input, actual output and review notes.
\item Export the weekly artifacts: conversion diagram; regression suite; code diff.
\end{itemize}
\textbf{Knowledge check.} Which security assumption was false? What exact precondition turns the programming error into a security vulnerability? What evidence would refute the finding?

\textbf{Scoring (10 points).} 3 root-cause/trace, 2 adversary conditions, 2 correction, 2 executable before/after evidence, 1 limits and clarity.

\textbf{Exit criteria.} A reviewer can reconstruct the flaw and confirm the relevant invariant now holds using the supplied tests; uncertain severity and environmental dependencies are recorded.

\textbf{Common pitfall.} Assuming all Windows integer-sized types have the same width.

\textbf{Next dependency.} Transfer the acquired audit pattern into the next original textbook chapter or the capstone; do not advance if trace, fix or evidence requirements are missing.

\subsection{Week 07 | W3 | Ownership, handles and failure returns}
\textbf{Prerequisite:} F1--F4; preceding specialized audit lab\hfill\textbf{Time:} 2h instruction + 2h reading + 4h lab

\textbf{Source reading [textbook-grounded]:} Ch. 7, Program Building Blocks, pp. 297--386; focus: Return Value Testing; Allocation Functions.

\textbf{Contemporary extension [not from textbook]:} RAII wrappers, SAL annotations and C++ static analysis

\textbf{Competency objectives (analyze/evaluate/create):}
\begin{itemize}[leftmargin=*,itemsep=2pt]
\item Diagram the input, privileged decision or state transition implicated in ownership, handles and failure returns.
\item Audit the nominated source operations and explain security conditions for review a routine that opens a HANDLE and exits on multiple branches.
\item Design a correction with both a positive test and an adversarial negative case; defend the residual-risk conclusion.
\end{itemize}
\textbf{Lecture and instructor demonstration.} Teach the relevant mechanism from Chapter 7, including the interaction between input representation, trust relationships, control flow and impacts. Instructor narrates a source-level trace, identifies a missing invariant, predicts a failing test, and demonstrates a corrected control. Link the older mechanism to RAII wrappers, SAL annotations and C++ static analysis, making version differences explicit.

\begin{insight}\textbf{Lab W3:} Review a routine that opens a HANDLE and exits on multiple branches.\\\textbf{Environment:} isolated Windows evaluation VM; Visual Studio Build Tools, WinDbg, PowerShell\\\textbf{Security assumption:} The sample is local/owned and untrusted inputs cannot be assumed canonical, authorized, or correctly sized.\end{insight}
\textbf{Starter code review miniature (new instructional example; not textbook text):}
\begin{lstlisting}[language=C]
/* Ownership review: second free is erroneous. */
char *p = malloc(24);
if (!p) return;
free(p);
/* free(p);  // double-free; remove */
/* Also review all early-return cleanup paths. */
\end{lstlisting}
\textbf{Laboratory execution procedure}
\begin{itemize}[leftmargin=*,itemsep=2pt]
\item Prepare a fresh working copy. Save the miniature as a discussion fixture, or map its flaw to the provided local application. Record baseline behavior and tool versions.
\item Formulate the security invariant and draw a source-to-sink or state-transition trace. List exactly which inputs, principals or concurrent transitions are under adversarial control.
\item Run controlled verification: Prove handle close-on-all-paths through tests and static review; check error-return usage.
\item Patch the underlying defect using a guard, authorization check, atomic operation, API contract or safe construction appropriate to the platform; note defense-in-depth controls separately.
\item Rerun the original case plus a valid control, an invalid case and a relevant boundary case. Save executable test input, actual output and review notes.
\item Export the weekly artifacts: path coverage table; RAII rewrite; static report.
\end{itemize}
\textbf{Knowledge check.} Which security assumption was false? What exact precondition turns the programming error into a security vulnerability? What evidence would refute the finding?

\textbf{Scoring (10 points).} 3 root-cause/trace, 2 adversary conditions, 2 correction, 2 executable before/after evidence, 1 limits and clarity.

\textbf{Exit criteria.} A reviewer can reconstruct the flaw and confirm the relevant invariant now holds using the supplied tests; uncertain severity and environmental dependencies are recorded.

\textbf{Common pitfall.} Treating INVALID\_HANDLE\_VALUE and NULL as interchangeable everywhere.

\textbf{Next dependency.} Transfer the acquired audit pattern into the next original textbook chapter or the capstone; do not advance if trace, fix or evidence requirements are missing.

\subsection{Week 08 | W4 | Unicode, filenames and dynamic execution}
\textbf{Prerequisite:} F1--F4; preceding specialized audit lab\hfill\textbf{Time:} 2h instruction + 2h reading + 4h lab

\textbf{Source reading [textbook-grounded]:} Ch. 8, Strings and Metacharacters, pp. 387--458; focus: Windows Unicode Functions; Path Metacharacters.

\textbf{Contemporary extension [not from textbook]:} UTF-16 and safe process creation

\textbf{Competency objectives (analyze/evaluate/create):}
\begin{itemize}[leftmargin=*,itemsep=2pt]
\item Diagram the input, privileged decision or state transition implicated in unicode, filenames and dynamic execution.
\item Audit the nominated source operations and explain security conditions for examine Windows filename normalization and argument quoting in a test helper.
\item Design a correction with both a positive test and an adversarial negative case; defend the residual-risk conclusion.
\end{itemize}
\textbf{Lecture and instructor demonstration.} Teach the relevant mechanism from Chapter 8, including the interaction between input representation, trust relationships, control flow and impacts. Instructor narrates a source-level trace, identifies a missing invariant, predicts a failing test, and demonstrates a corrected control. Link the older mechanism to UTF-16 and safe process creation, making version differences explicit.

\begin{insight}\textbf{Lab W4:} Examine Windows filename normalization and argument quoting in a test helper.\\\textbf{Environment:} isolated Windows evaluation VM; Visual Studio Build Tools, WinDbg, PowerShell\\\textbf{Security assumption:} The sample is local/owned and untrusted inputs cannot be assumed canonical, authorized, or correctly sized.\end{insight}
\textbf{Starter code review miniature (new instructional example; not textbook text):}
\begin{lstlisting}[language=Python]
# Unsafe when the filename is attacker-controlled:
subprocess.run("cat " + filename, shell=True)
# Safer only for an authorized, validated file target:
subprocess.run(["cat", "--", filename], check=True)
# Also review path containment and effective permissions.
\end{lstlisting}
\textbf{Laboratory execution procedure}
\begin{itemize}[leftmargin=*,itemsep=2pt]
\item Prepare a fresh working copy. Save the miniature as a discussion fixture, or map its flaw to the provided local application. Record baseline behavior and tool versions.
\item Formulate the security invariant and draw a source-to-sink or state-transition trace. List exactly which inputs, principals or concurrent transitions are under adversarial control.
\item Run controlled verification: Use spaces and mixed Unicode paths; compare raw command string to separated executable/arguments.
\item Patch the underlying defect using a guard, authorization check, atomic operation, API contract or safe construction appropriate to the platform; note defense-in-depth controls separately.
\item Rerun the original case plus a valid control, an invalid case and a relevant boundary case. Save executable test input, actual output and review notes.
\item Export the weekly artifacts: string transformation map; patched helper; tests.
\end{itemize}
\textbf{Knowledge check.} Which security assumption was false? What exact precondition turns the programming error into a security vulnerability? What evidence would refute the finding?

\textbf{Scoring (10 points).} 3 root-cause/trace, 2 adversary conditions, 2 correction, 2 executable before/after evidence, 1 limits and clarity.

\textbf{Exit criteria.} A reviewer can reconstruct the flaw and confirm the relevant invariant now holds using the supplied tests; uncertain severity and environmental dependencies are recorded.

\textbf{Common pitfall.} Equating text escaping with reliable argument separation.

\textbf{Next dependency.} Transfer the acquired audit pattern into the next original textbook chapter or the capstone; do not advance if trace, fix or evidence requirements are missing.

\begin{historynote}\textbf{Checkpoint:} Midterm practical: source audit under time limits, marked for evidence rather than tool-alert count.\end{historynote}
\subsection{Week 09 | W5 | Tokens and object access controls}
\textbf{Prerequisite:} F1--F4; preceding specialized audit lab\hfill\textbf{Time:} 2h instruction + 2h reading + 4h lab

\textbf{Source reading [textbook-grounded]:} Ch. 11, Windows I: Objects and the File System, pp. 625--684; focus: Access Tokens; Security Descriptors.

\textbf{Contemporary extension [not from textbook]:} integrity levels, service SIDs and least privilege

\textbf{Competency objectives (analyze/evaluate/create):}
\begin{itemize}[leftmargin=*,itemsep=2pt]
\item Diagram the input, privileged decision or state transition implicated in tokens and object access controls.
\item Audit the nominated source operations and explain security conditions for inspect a Windows service object illustrative DACL and token permissions.
\item Design a correction with both a positive test and an adversarial negative case; defend the residual-risk conclusion.
\end{itemize}
\textbf{Lecture and instructor demonstration.} Teach the relevant mechanism from Chapter 11, including the interaction between input representation, trust relationships, control flow and impacts. Instructor narrates a source-level trace, identifies a missing invariant, predicts a failing test, and demonstrates a corrected control. Link the older mechanism to integrity levels, service SIDs and least privilege, making version differences explicit.

\begin{insight}\textbf{Lab W5:} Inspect a Windows service object illustrative DACL and token permissions.\\\textbf{Environment:} isolated Windows evaluation VM; Visual Studio Build Tools, WinDbg, PowerShell\\\textbf{Security assumption:} The sample is local/owned and untrusted inputs cannot be assumed canonical, authorized, or correctly sized.\end{insight}
\textbf{Starter code review miniature (new instructional example; not textbook text):}
\begin{lstlisting}[language=C]
/* Windows API: check return and effective ACLs. */
HANDLE h = CreateFileW(path, GENERIC_READ, FILE_SHARE_READ,
    NULL, OPEN_EXISTING, FILE_ATTRIBUTE_NORMAL, NULL);
if (h == INVALID_HANDLE_VALUE) return FALSE;
CloseHandle(h);
/* Caller access must be separately authorized. */
\end{lstlisting}
\textbf{Laboratory execution procedure}
\begin{itemize}[leftmargin=*,itemsep=2pt]
\item Prepare a fresh working copy. Save the miniature as a discussion fixture, or map its flaw to the provided local application. Record baseline behavior and tool versions.
\item Formulate the security invariant and draw a source-to-sink or state-transition trace. List exactly which inputs, principals or concurrent transitions are under adversarial control.
\item Run controlled verification: Check allowed and denied principals with an isolated Windows VM; avoid modifying host permissions.
\item Patch the underlying defect using a guard, authorization check, atomic operation, API contract or safe construction appropriate to the platform; note defense-in-depth controls separately.
\item Rerun the original case plus a valid control, an invalid case and a relevant boundary case. Save executable test input, actual output and review notes.
\item Export the weekly artifacts: ACL matrix; token reasoning; PowerShell verification output.
\end{itemize}
\textbf{Knowledge check.} Which security assumption was false? What exact precondition turns the programming error into a security vulnerability? What evidence would refute the finding?

\textbf{Scoring (10 points).} 3 root-cause/trace, 2 adversary conditions, 2 correction, 2 executable before/after evidence, 1 limits and clarity.

\textbf{Exit criteria.} A reviewer can reconstruct the flaw and confirm the relevant invariant now holds using the supplied tests; uncertain severity and environmental dependencies are recorded.

\textbf{Common pitfall.} Conflating token privilege presence with enabled and authorized access.

\textbf{Next dependency.} Transfer the acquired audit pattern into the next original textbook chapter or the capstone; do not advance if trace, fix or evidence requirements are missing.

\subsection{Week 10 | W6 | Files, registry and DLL loading}
\textbf{Prerequisite:} F1--F4; preceding specialized audit lab\hfill\textbf{Time:} 2h instruction + 2h reading + 4h lab

\textbf{Source reading [textbook-grounded]:} Ch. 11, Windows I: Objects and the File System, pp. 625--684; focus: DLL Loading; File Access; The Registry.

\textbf{Contemporary extension [not from textbook]:} safe search directories, ACL inheritance and service settings

\textbf{Competency objectives (analyze/evaluate/create):}
\begin{itemize}[leftmargin=*,itemsep=2pt]
\item Diagram the input, privileged decision or state transition implicated in files, registry and dll loading.
\item Audit the nominated source operations and explain security conditions for review insecure DLL resolution and overbroad registry-key ACLs.
\item Design a correction with both a positive test and an adversarial negative case; defend the residual-risk conclusion.
\end{itemize}
\textbf{Lecture and instructor demonstration.} Teach the relevant mechanism from Chapter 11, including the interaction between input representation, trust relationships, control flow and impacts. Instructor narrates a source-level trace, identifies a missing invariant, predicts a failing test, and demonstrates a corrected control. Link the older mechanism to safe search directories, ACL inheritance and service settings, making version differences explicit.

\begin{insight}\textbf{Lab W6:} Review insecure DLL resolution and overbroad registry-key ACLs.\\\textbf{Environment:} isolated Windows evaluation VM; Visual Studio Build Tools, WinDbg, PowerShell\\\textbf{Security assumption:} The sample is local/owned and untrusted inputs cannot be assumed canonical, authorized, or correctly sized.\end{insight}
\textbf{Starter code review miniature (new instructional example; not textbook text):}
\begin{lstlisting}[language=C]
/* Windows API: check return and effective ACLs. */
HANDLE h = CreateFileW(path, GENERIC_READ, FILE_SHARE_READ,
    NULL, OPEN_EXISTING, FILE_ATTRIBUTE_NORMAL, NULL);
if (h == INVALID_HANDLE_VALUE) return FALSE;
CloseHandle(h);
/* Caller access must be separately authorized. */
\end{lstlisting}
\textbf{Laboratory execution procedure}
\begin{itemize}[leftmargin=*,itemsep=2pt]
\item Prepare a fresh working copy. Save the miniature as a discussion fixture, or map its flaw to the provided local application. Record baseline behavior and tool versions.
\item Formulate the security invariant and draw a source-to-sink or state-transition trace. List exactly which inputs, principals or concurrent transitions are under adversarial control.
\item Run controlled verification: Validate search-order assumptions with a benign test DLL only in an isolated VM; remediate absolute loading.
\item Patch the underlying defect using a guard, authorization check, atomic operation, API contract or safe construction appropriate to the platform; note defense-in-depth controls separately.
\item Rerun the original case plus a valid control, an invalid case and a relevant boundary case. Save executable test input, actual output and review notes.
\item Export the weekly artifacts: load trace; ACL before/after; mitigation evidence.
\end{itemize}
\textbf{Knowledge check.} Which security assumption was false? What exact precondition turns the programming error into a security vulnerability? What evidence would refute the finding?

\textbf{Scoring (10 points).} 3 root-cause/trace, 2 adversary conditions, 2 correction, 2 executable before/after evidence, 1 limits and clarity.

\textbf{Exit criteria.} A reviewer can reconstruct the flaw and confirm the relevant invariant now holds using the supplied tests; uncertain severity and environmental dependencies are recorded.

\textbf{Common pitfall.} Accepting writable directories in a privileged DLL search path.

\textbf{Next dependency.} Transfer the acquired audit pattern into the next original textbook chapter or the capstone; do not advance if trace, fix or evidence requirements are missing.

\subsection{Week 11 | W7 | Named pipes and impersonation}
\textbf{Prerequisite:} F1--F4; preceding specialized audit lab\hfill\textbf{Time:} 2h instruction + 2h reading + 4h lab

\textbf{Source reading [textbook-grounded]:} Ch. 12, Windows II: Interprocess Communication, pp. 685--754; focus: Pipes; Impersonation in Pipes.

\textbf{Contemporary extension [not from textbook]:} pipe ACLs and authenticated client security context

\textbf{Competency objectives (analyze/evaluate/create):}
\begin{itemize}[leftmargin=*,itemsep=2pt]
\item Diagram the input, privileged decision or state transition implicated in named pipes and impersonation.
\item Audit the nominated source operations and explain security conditions for audit a sample service that trusts all named-pipe clients.
\item Design a correction with both a positive test and an adversarial negative case; defend the residual-risk conclusion.
\end{itemize}
\textbf{Lecture and instructor demonstration.} Teach the relevant mechanism from Chapter 12, including the interaction between input representation, trust relationships, control flow and impacts. Instructor narrates a source-level trace, identifies a missing invariant, predicts a failing test, and demonstrates a corrected control. Link the older mechanism to pipe ACLs and authenticated client security context, making version differences explicit.

\begin{insight}\textbf{Lab W7:} Audit a sample service that trusts all named-pipe clients.\\\textbf{Environment:} isolated Windows evaluation VM; Visual Studio Build Tools, WinDbg, PowerShell\\\textbf{Security assumption:} The sample is local/owned and untrusted inputs cannot be assumed canonical, authorized, or correctly sized.\end{insight}
\textbf{Starter code review miniature (new instructional example; not textbook text):}
\begin{lstlisting}[language=C]
/* Windows named-pipe: a connection is not authorization. */
BOOL connected = ConnectNamedPipe(hPipe, NULL);
/* Review ERROR_PIPE_CONNECTED and API error cases. */
/* Check client identity, pipe DACL and method permissions. */
\end{lstlisting}
\textbf{Laboratory execution procedure}
\begin{itemize}[leftmargin=*,itemsep=2pt]
\item Prepare a fresh working copy. Save the miniature as a discussion fixture, or map its flaw to the provided local application. Record baseline behavior and tool versions.
\item Formulate the security invariant and draw a source-to-sink or state-transition trace. List exactly which inputs, principals or concurrent transitions are under adversarial control.
\item Run controlled verification: Specify and test a restrictive pipe security descriptor and a client-identity check.
\item Patch the underlying defect using a guard, authorization check, atomic operation, API contract or safe construction appropriate to the platform; note defense-in-depth controls separately.
\item Rerun the original case plus a valid control, an invalid case and a relevant boundary case. Save executable test input, actual output and review notes.
\item Export the weekly artifacts: pipe trust diagram; fixed checks; access test.
\end{itemize}
\textbf{Knowledge check.} Which security assumption was false? What exact precondition turns the programming error into a security vulnerability? What evidence would refute the finding?

\textbf{Scoring (10 points).} 3 root-cause/trace, 2 adversary conditions, 2 correction, 2 executable before/after evidence, 1 limits and clarity.

\textbf{Exit criteria.} A reviewer can reconstruct the flaw and confirm the relevant invariant now holds using the supplied tests; uncertain severity and environmental dependencies are recorded.

\textbf{Common pitfall.} Assuming a local connection is inherently trusted.

\textbf{Next dependency.} Transfer the acquired audit pattern into the next original textbook chapter or the capstone; do not advance if trace, fix or evidence requirements are missing.

\subsection{Week 12 | W8 | RPC, COM and service boundaries}
\textbf{Prerequisite:} F1--F4; preceding specialized audit lab\hfill\textbf{Time:} 2h instruction + 2h reading + 4h lab

\textbf{Source reading [textbook-grounded]:} Ch. 12, Windows II: Interprocess Communication, pp. 685--754; focus: Auditing RPC Applications; COM; DCOM Access Controls.

\textbf{Contemporary extension [not from textbook]:} RPC authorization callbacks and COM identity

\textbf{Competency objectives (analyze/evaluate/create):}
\begin{itemize}[leftmargin=*,itemsep=2pt]
\item Diagram the input, privileged decision or state transition implicated in rpc, com and service boundaries.
\item Audit the nominated source operations and explain security conditions for trace a privileged service method exposed through RPC/COM in a provided interface definition.
\item Design a correction with both a positive test and an adversarial negative case; defend the residual-risk conclusion.
\end{itemize}
\textbf{Lecture and instructor demonstration.} Teach the relevant mechanism from Chapter 12, including the interaction between input representation, trust relationships, control flow and impacts. Instructor narrates a source-level trace, identifies a missing invariant, predicts a failing test, and demonstrates a corrected control. Link the older mechanism to RPC authorization callbacks and COM identity, making version differences explicit.

\begin{insight}\textbf{Lab W8:} Trace a privileged service method exposed through RPC/COM in a provided interface definition.\\\textbf{Environment:} isolated Windows evaluation VM; Visual Studio Build Tools, WinDbg, PowerShell\\\textbf{Security assumption:} The sample is local/owned and untrusted inputs cannot be assumed canonical, authorized, or correctly sized.\end{insight}
\textbf{Starter code review miniature (new instructional example; not textbook text):}
\begin{lstlisting}[language=C]
/* Windows named-pipe: a connection is not authorization. */
BOOL connected = ConnectNamedPipe(hPipe, NULL);
/* Review ERROR_PIPE_CONNECTED and API error cases. */
/* Check client identity, pipe DACL and method permissions. */
\end{lstlisting}
\textbf{Laboratory execution procedure}
\begin{itemize}[leftmargin=*,itemsep=2pt]
\item Prepare a fresh working copy. Save the miniature as a discussion fixture, or map its flaw to the provided local application. Record baseline behavior and tool versions.
\item Formulate the security invariant and draw a source-to-sink or state-transition trace. List exactly which inputs, principals or concurrent transitions are under adversarial control.
\item Run controlled verification: Check per-operation authorization and interface-level exposure using benign calls.
\item Patch the underlying defect using a guard, authorization check, atomic operation, API contract or safe construction appropriate to the platform; note defense-in-depth controls separately.
\item Rerun the original case plus a valid control, an invalid case and a relevant boundary case. Save executable test input, actual output and review notes.
\item Export the weekly artifacts: call graph; authorization matrix; test protocol.
\end{itemize}
\textbf{Knowledge check.} Which security assumption was false? What exact precondition turns the programming error into a security vulnerability? What evidence would refute the finding?

\textbf{Scoring (10 points).} 3 root-cause/trace, 2 adversary conditions, 2 correction, 2 executable before/after evidence, 1 limits and clarity.

\textbf{Exit criteria.} A reviewer can reconstruct the flaw and confirm the relevant invariant now holds using the supplied tests; uncertain severity and environmental dependencies are recorded.

\textbf{Common pitfall.} Trusting transport authentication as sufficient per-method authorization.

\textbf{Next dependency.} Transfer the acquired audit pattern into the next original textbook chapter or the capstone; do not advance if trace, fix or evidence requirements are missing.

\subsection{Week 13 | W9 | Windows thread safety and races}
\textbf{Prerequisite:} F1--F4; preceding specialized audit lab\hfill\textbf{Time:} 2h instruction + 2h reading + 4h lab

\textbf{Source reading [textbook-grounded]:} Ch. 13, Synchronization and State, pp. 755--828; focus: Windows Process Synchronization; Threading Vulnerabilities.

\textbf{Contemporary extension [not from textbook]:} SRW locks, interlocked operations and concurrency testing

\textbf{Competency objectives (analyze/evaluate/create):}
\begin{itemize}[leftmargin=*,itemsep=2pt]
\item Diagram the input, privileged decision or state transition implicated in windows thread safety and races.
\item Audit the nominated source operations and explain security conditions for review two worker threads updating shared authorization state.
\item Design a correction with both a positive test and an adversarial negative case; defend the residual-risk conclusion.
\end{itemize}
\textbf{Lecture and instructor demonstration.} Teach the relevant mechanism from Chapter 13, including the interaction between input representation, trust relationships, control flow and impacts. Instructor narrates a source-level trace, identifies a missing invariant, predicts a failing test, and demonstrates a corrected control. Link the older mechanism to SRW locks, interlocked operations and concurrency testing, making version differences explicit.

\begin{insight}\textbf{Lab W9:} Review two worker threads updating shared authorization state.\\\textbf{Environment:} isolated Windows evaluation VM; Visual Studio Build Tools, WinDbg, PowerShell\\\textbf{Security assumption:} The sample is local/owned and untrusted inputs cannot be assumed canonical, authorized, or correctly sized.\end{insight}
\textbf{Starter code review miniature (new instructional example; not textbook text):}
\begin{lstlisting}[language=C]
/* This read-modify-write is not a synchronized quota. */
int quota = 10;
void spend(void) { if (quota > 0) quota--; }
/* Concurrent callers require an atomic state transition
   or a correctly scoped lock. */
\end{lstlisting}
\textbf{Laboratory execution procedure}
\begin{itemize}[leftmargin=*,itemsep=2pt]
\item Prepare a fresh working copy. Save the miniature as a discussion fixture, or map its flaw to the provided local application. Record baseline behavior and tool versions.
\item Formulate the security invariant and draw a source-to-sink or state-transition trace. List exactly which inputs, principals or concurrent transitions are under adversarial control.
\item Run controlled verification: Produce a schedule demonstrating stale or inconsistent state and validate synchronization correction.
\item Patch the underlying defect using a guard, authorization check, atomic operation, API contract or safe construction appropriate to the platform; note defense-in-depth controls separately.
\item Rerun the original case plus a valid control, an invalid case and a relevant boundary case. Save executable test input, actual output and review notes.
\item Export the weekly artifacts: interleaving chart; lock protocol; tests.
\end{itemize}
\textbf{Knowledge check.} Which security assumption was false? What exact precondition turns the programming error into a security vulnerability? What evidence would refute the finding?

\textbf{Scoring (10 points).} 3 root-cause/trace, 2 adversary conditions, 2 correction, 2 executable before/after evidence, 1 limits and clarity.

\textbf{Exit criteria.} A reviewer can reconstruct the flaw and confirm the relevant invariant now holds using the supplied tests; uncertain severity and environmental dependencies are recorded.

\textbf{Common pitfall.} Protecting writes while leaving unsynchronized reads.

\textbf{Next dependency.} Transfer the acquired audit pattern into the next original textbook chapter or the capstone; do not advance if trace, fix or evidence requirements are missing.

\subsection{Week 14 | W10 | Windows native audit integration}
\textbf{Prerequisite:} F1--F4; preceding specialized audit lab\hfill\textbf{Time:} 2h instruction + 2h reading + 4h lab

\textbf{Source reading [textbook-grounded]:} Ch. 13, Synchronization and State, pp. 755--828; focus: Reentrancy; Race Conditions; Threading Vulnerabilities.

\textbf{Contemporary extension [not from textbook]:} WinDbg, ETW and CodeQL C/C++

\textbf{Competency objectives (analyze/evaluate/create):}
\begin{itemize}[leftmargin=*,itemsep=2pt]
\item Diagram the input, privileged decision or state transition implicated in windows native audit integration.
\item Audit the nominated source operations and explain security conditions for combine findings from token, IPC and thread reviews in a service assessment.
\item Design a correction with both a positive test and an adversarial negative case; defend the residual-risk conclusion.
\end{itemize}
\textbf{Lecture and instructor demonstration.} Teach the relevant mechanism from Chapter 13, including the interaction between input representation, trust relationships, control flow and impacts. Instructor narrates a source-level trace, identifies a missing invariant, predicts a failing test, and demonstrates a corrected control. Link the older mechanism to WinDbg, ETW and CodeQL C/C++, making version differences explicit.

\begin{insight}\textbf{Lab W10:} Combine findings from token, IPC and thread reviews in a service assessment.\\\textbf{Environment:} isolated Windows evaluation VM; Visual Studio Build Tools, WinDbg, PowerShell\\\textbf{Security assumption:} The sample is local/owned and untrusted inputs cannot be assumed canonical, authorized, or correctly sized.\end{insight}
\textbf{Starter code review miniature (new instructional example; not textbook text):}
\begin{lstlisting}[language=C]
/* This read-modify-write is not a synchronized quota. */
int quota = 10;
void spend(void) { if (quota > 0) quota--; }
/* Concurrent callers require an atomic state transition
   or a correctly scoped lock. */
\end{lstlisting}
\textbf{Laboratory execution procedure}
\begin{itemize}[leftmargin=*,itemsep=2pt]
\item Prepare a fresh working copy. Save the miniature as a discussion fixture, or map its flaw to the provided local application. Record baseline behavior and tool versions.
\item Formulate the security invariant and draw a source-to-sink or state-transition trace. List exactly which inputs, principals or concurrent transitions are under adversarial control.
\item Run controlled verification: Reproduce an allowed/denied state change in a controlled VM and rerun after hardening.
\item Patch the underlying defect using a guard, authorization check, atomic operation, API contract or safe construction appropriate to the platform; note defense-in-depth controls separately.
\item Rerun the original case plus a valid control, an invalid case and a relevant boundary case. Save executable test input, actual output and review notes.
\item Export the weekly artifacts: consolidated report; test runs; residual-risk memo.
\end{itemize}
\textbf{Knowledge check.} Which security assumption was false? What exact precondition turns the programming error into a security vulnerability? What evidence would refute the finding?

\textbf{Scoring (10 points).} 3 root-cause/trace, 2 adversary conditions, 2 correction, 2 executable before/after evidence, 1 limits and clarity.

\textbf{Exit criteria.} A reviewer can reconstruct the flaw and confirm the relevant invariant now holds using the supplied tests; uncertain severity and environmental dependencies are recorded.

\textbf{Common pitfall.} Over-claiming WinDbg traces as proof of non-existence of a vulnerability.

\textbf{Next dependency.} Transfer the acquired audit pattern into the next original textbook chapter or the capstone; do not advance if trace, fix or evidence requirements are missing.

\begin{historynote}\textbf{Checkpoint:} Specialization signoff: independently audit a comparable unseen fixture before capstone.\end{historynote}
\subsection{Week 15 | C1 | Capstone: audit and evidence}
\textbf{Prerequisite:} all assessed specialization competencies\hfill\textbf{Time:} 2h instruction + 2h reading + 4h lab

\textbf{Source reading [textbook-grounded]:} Ch. 4, Application Review Process, pp. 91--166; focus: Plan--Work--Reflect; Reporting and Remediation Support.

\textbf{Contemporary extension [not from textbook]:} SARIF evidence, pull-request review, regression tests and residual risk

\textbf{Competency objectives (analyze/evaluate/create):}
\begin{itemize}[leftmargin=*,itemsep=2pt]
\item Diagram the input, privileged decision or state transition implicated in capstone: audit and evidence.
\item Audit the nominated source operations and explain security conditions for apply the complete assessment method to the specialization capstone fixture, with scoped objectives and defensible vulnerability claims.
\item Design a correction with both a positive test and an adversarial negative case; defend the residual-risk conclusion.
\end{itemize}
\textbf{Lecture and instructor demonstration.} Teach the relevant mechanism from Chapter 4, including the interaction between input representation, trust relationships, control flow and impacts. Instructor narrates a source-level trace, identifies a missing invariant, predicts a failing test, and demonstrates a corrected control. Link the older mechanism to SARIF evidence, pull-request review, regression tests and residual risk, making version differences explicit.

\begin{insight}\textbf{Lab C1:} Apply the complete assessment method to the specialization capstone fixture, with scoped objectives and defensible vulnerability claims.\\\textbf{Environment:} isolated Windows evaluation VM; Visual Studio Build Tools, WinDbg, PowerShell\\\textbf{Security assumption:} The sample is local/owned and untrusted inputs cannot be assumed canonical, authorized, or correctly sized.\end{insight}
\textbf{Starter code review miniature (new instructional example; not textbook text):}
\begin{lstlisting}[language=Python]
def triage(source, sink, guard):
    # A finding requires more than matching a dangerous call.
    return {"source": source, "sink": sink,
            "guard": guard, "reachable": None,
            "validated": False}
# Auditor must investigate real guards and execution paths.
\end{lstlisting}
\textbf{Laboratory execution procedure}
\begin{itemize}[leftmargin=*,itemsep=2pt]
\item Prepare a fresh working copy. Save the miniature as a discussion fixture, or map its flaw to the provided local application. Record baseline behavior and tool versions.
\item Formulate the security invariant and draw a source-to-sink or state-transition trace. List exactly which inputs, principals or concurrent transitions are under adversarial control.
\item Run controlled verification: In week 15, establish reachability with controlled tests and trace evidence; do not confuse conjecture with confirmed findings.
\item Patch the underlying defect using a guard, authorization check, atomic operation, API contract or safe construction appropriate to the platform; note defense-in-depth controls separately.
\item Rerun the original case plus a valid control, an invalid case and a relevant boundary case. Save executable test input, actual output and review notes.
\item Export the weekly artifacts: scope; architecture; threat model; annotated findings; controlled validation logs.
\end{itemize}
\textbf{Knowledge check.} Which security assumption was false? What exact precondition turns the programming error into a security vulnerability? What evidence would refute the finding?

\textbf{Scoring (10 points).} 3 root-cause/trace, 2 adversary conditions, 2 correction, 2 executable before/after evidence, 1 limits and clarity.

\textbf{Exit criteria.} A reviewer can reconstruct the flaw and confirm the relevant invariant now holds using the supplied tests; uncertain severity and environmental dependencies are recorded.

\textbf{Common pitfall.} Incomplete scoping or unsupported exploitability assumptions.

\textbf{Next dependency.} Transfer the acquired audit pattern into the next original textbook chapter or the capstone; do not advance if trace, fix or evidence requirements are missing.

\subsection{Week 16 | C2 | Capstone: remediation and assurance}
\textbf{Prerequisite:} all assessed specialization competencies\hfill\textbf{Time:} 2h instruction + 2h reading + 4h lab

\textbf{Source reading [textbook-grounded]:} Ch. 4, Application Review Process, pp. 91--166; focus: Plan--Work--Reflect; Reporting and Remediation Support.

\textbf{Contemporary extension [not from textbook]:} SARIF evidence, pull-request review, regression tests and residual risk

\textbf{Competency objectives (analyze/evaluate/create):}
\begin{itemize}[leftmargin=*,itemsep=2pt]
\item Diagram the input, privileged decision or state transition implicated in capstone: remediation and assurance.
\item Audit the nominated source operations and explain security conditions for implement and test fixes; present a complete technical report, evidence appendix and residual-risk briefing.
\item Design a correction with both a positive test and an adversarial negative case; defend the residual-risk conclusion.
\end{itemize}
\textbf{Lecture and instructor demonstration.} Teach the relevant mechanism from Chapter 4, including the interaction between input representation, trust relationships, control flow and impacts. Instructor narrates a source-level trace, identifies a missing invariant, predicts a failing test, and demonstrates a corrected control. Link the older mechanism to SARIF evidence, pull-request review, regression tests and residual risk, making version differences explicit.

\begin{insight}\textbf{Lab C2:} Implement and test fixes; present a complete technical report, evidence appendix and residual-risk briefing.\\\textbf{Environment:} isolated Windows evaluation VM; Visual Studio Build Tools, WinDbg, PowerShell\\\textbf{Security assumption:} The sample is local/owned and untrusted inputs cannot be assumed canonical, authorized, or correctly sized.\end{insight}
\textbf{Starter code review miniature (new instructional example; not textbook text):}
\begin{lstlisting}[language=Python]
def triage(source, sink, guard):
    # A finding requires more than matching a dangerous call.
    return {"source": source, "sink": sink,
            "guard": guard, "reachable": None,
            "validated": False}
# Auditor must investigate real guards and execution paths.
\end{lstlisting}
\textbf{Laboratory execution procedure}
\begin{itemize}[leftmargin=*,itemsep=2pt]
\item Prepare a fresh working copy. Save the miniature as a discussion fixture, or map its flaw to the provided local application. Record baseline behavior and tool versions.
\item Formulate the security invariant and draw a source-to-sink or state-transition trace. List exactly which inputs, principals or concurrent transitions are under adversarial control.
\item Run controlled verification: Rerun the original reproducer plus negative controls; obtain independent reviewer signoff.
\item Patch the underlying defect using a guard, authorization check, atomic operation, API contract or safe construction appropriate to the platform; note defense-in-depth controls separately.
\item Rerun the original case plus a valid control, an invalid case and a relevant boundary case. Save executable test input, actual output and review notes.
\item Export the weekly artifacts: patchset; test evidence; SARIF/triage record; final written report and presentation.
\end{itemize}
\textbf{Knowledge check.} Which security assumption was false? What exact precondition turns the programming error into a security vulnerability? What evidence would refute the finding?

\textbf{Scoring (10 points).} 3 root-cause/trace, 2 adversary conditions, 2 correction, 2 executable before/after evidence, 1 limits and clarity.

\textbf{Exit criteria.} A reviewer can reconstruct the flaw and confirm the relevant invariant now holds using the supplied tests; uncertain severity and environmental dependencies are recorded.

\textbf{Common pitfall.} Declaring remediation complete without exercising the original failure mode.

\textbf{Next dependency.} Transfer the acquired audit pattern into the next original textbook chapter or the capstone; do not advance if trace, fix or evidence requirements are missing.

\begin{historynote}\textbf{Checkpoint:} Professional portfolio defense: peer reviewer challenges assumptions, reproducers, fixes and residual risk.\end{historynote}
\section{Laboratory catalog and traceable evidence}
\begingroup\footnotesize\setlength\tabcolsep{3pt}
\begin{longtable}{>{\raggedright\arraybackslash}p{\dimexpr0.100\textwidth-6pt\relax}>{\raggedright\arraybackslash}p{\dimexpr0.480\textwidth-6pt\relax}>{\raggedright\arraybackslash}p{\dimexpr0.420\textwidth-6pt\relax}}\toprule
\textbf{Lab} & \textbf{Source / original theme} & \textbf{Assessment work product}\\\midrule\endfirsthead
\toprule \textbf{Lab} & \textbf{Source / original theme} & \textbf{Assessment work product}\\\midrule\endhead
F1 & Ch. 1: Vulnerabilities; Security Policies; Classifying Vulnerabilities & classification.csv; 1-page policy assumptions; issue evidence notes\\\midrule
F2 & Ch. 2: Trust Relationships; Threat Modeling; Architecture Modeling & DFD; threat register; ranked audit entry points\\\midrule
F3 & Ch. 3: Exposure; Insecure Defaults; Protective Measures & exposure matrix; hardened configuration diff; verification log\\\midrule
F4 & Ch. 4: Code-Auditing Strategies; Internal Flow Analysis; Reporting & annotated flow trace; triage record; remediation PR with regression tests\\\midrule
W1 & Ch. 5: Buffer Overflows; Protection Mechanisms & memory map; corrected code; validation cases\\\midrule
W2 & Ch. 6: Signed/Unsigned Conversions; Truncation & conversion diagram; regression suite; code diff\\\midrule
W3 & Ch. 7: Return Value Testing; Allocation Functions & path coverage table; RAII rewrite; static report\\\midrule
W4 & Ch. 8: Windows Unicode Functions; Path Metacharacters & string transformation map; patched helper; tests\\\midrule
W5 & Ch. 11: Access Tokens; Security Descriptors & ACL matrix; token reasoning; PowerShell verification output\\\midrule
W6 & Ch. 11: DLL Loading; File Access; The Registry & load trace; ACL before/after; mitigation evidence\\\midrule
W7 & Ch. 12: Pipes; Impersonation in Pipes & pipe trust diagram; fixed checks; access test\\\midrule
W8 & Ch. 12: Auditing RPC Applications; COM; DCOM Access Controls & call graph; authorization matrix; test protocol\\\midrule
W9 & Ch. 13: Windows Process Synchronization; Threading Vulnerabilities & interleaving chart; lock protocol; tests\\\midrule
W10 & Ch. 13: Reentrancy; Race Conditions; Threading Vulnerabilities & consolidated report; test runs; residual-risk memo\\\midrule
C1 & Ch. 4: Plan--Work--Reflect; Reporting and Remediation Support & scope; architecture; threat model; annotated findings; controlled validation logs\\\midrule
C2 & Ch. 4: Plan--Work--Reflect; Reporting and Remediation Support & patchset; test evidence; SARIF/triage record; final written report and presentation\\\midrule
\bottomrule\end{longtable}\endgroup
\section{Tooling, availability and laboratory environments}
All learners can complete the core manual analysis with a text editor, compiler/interpreter, test framework, and a disposable local environment. Paid services are optional, not mandatory prerequisites.

\begingroup\footnotesize\setlength\tabcolsep{3pt}
\begin{longtable}{>{\raggedright\arraybackslash}p{\dimexpr0.200\textwidth-6pt\relax}>{\raggedright\arraybackslash}p{\dimexpr0.330\textwidth-6pt\relax}>{\raggedright\arraybackslash}p{\dimexpr0.470\textwidth-6pt\relax}}\toprule
\textbf{Tool} & \textbf{Competency supported} & \textbf{Limitations / setup}\\\midrule\endfirsthead
\toprule \textbf{Tool} & \textbf{Competency supported} & \textbf{Limitations / setup}\\\midrule\endhead
CodeQL CLI / GitHub code scanning & Semantic flow and query analysis; supports C/C++, C\#, Go, Java/Kotlin, JS/TS, Python, Rust and other listed languages; not PHP & Review GitHub licensing, API availability, supported language/framework versions and build mode before use\\\midrule
Semgrep & Pattern/dataflow analysis over configured supported languages & Use documented community/commercial features; rules require manual validation\\\midrule
Git, editor and tests & Manual inspection, reproducibility and change review & Local version control; avoid committing secrets or personal data\\\midrule
GitHub Actions & CI validation and gate demonstrations & Platform limits, repository policy and network access may apply\\\midrule
Visual Studio Build Tools / WinDbg & Windows native builds, debug and investigation & Windows-only; some workloads/features require installation\\\midrule
PowerShell + Windows SDK & Inspect DACLs, tokens and service configuration & Use isolated VM; no changes to host policies\\\midrule
CodeQL C/C++/C\# & Cross-file semantic assessment & Verify database build coverage and license\\\midrule
\bottomrule\end{longtable}\endgroup
\subsection{Reproducibility contract}
\begin{itemize}[leftmargin=*,itemsep=2pt]
\item Record OS, compiler/interpreter version, dependencies, invocation and expected exit status.
\item Use instructor-owned fixture repositories at a pinned commit; keep a checksum or revision identifier.
\item Isolate each lab using a VM, rootless container, local loopback, or offline fixtures.
\item Capture both positive and negative test results, including outputs after remediation.
\item If native platform access is unavailable, grade a clearly labeled static-analysis alternative rather than claiming live verification.
\end{itemize}
\section{Assessment, proficiency and assurance}
Assessment weights are consistent across all tracks; successful completion requires at least 70/100 overall, plus no omitted security evidence in the final capstone. Competency mastery is a separate judgment: a passing multiple-choice score alone never certifies independent auditing ability.

\begingroup\footnotesize\setlength\tabcolsep{3pt}
\begin{longtable}{>{\raggedright\arraybackslash}p{\dimexpr0.220\textwidth-6pt\relax}>{\raggedright\arraybackslash}p{\dimexpr0.110\textwidth-6pt\relax}>{\raggedright\arraybackslash}p{\dimexpr0.670\textwidth-6pt\relax}}\toprule
\textbf{Component} & \textbf{Weight} & \textbf{Measurement}\\\midrule\endfirsthead
\toprule \textbf{Component} & \textbf{Weight} & \textbf{Measurement}\\\midrule\endhead
Weekly laboratories & 30\% & Lab notebook, executable negative tests, verifiable artifacts and remediation review\\\midrule
Source-code auditing & 20\% & Source/sink trace, execution conditions, control evaluation and evidence quality\\\midrule
Knowledge assessments & 15\% & Short scenarios, design arguments, false-positive differentiation\\\midrule
Midterm practical & 10\% & Timed code audit in week 8 with independent reproduction\\\midrule
Capstone & 25\% & Scoped independent analysis, corrections, tests and risk communication\\\midrule
\bottomrule\end{longtable}\endgroup
\subsection{Observable proficiency rubric}
\begingroup\footnotesize\setlength\tabcolsep{3pt}
\begin{longtable}{>{\raggedright\arraybackslash}p{\dimexpr0.200\textwidth-6pt\relax}>{\raggedright\arraybackslash}p{\dimexpr0.800\textwidth-6pt\relax}}\toprule
\textbf{Level} & \textbf{Evidence of competency}\\\midrule\endfirsthead
\toprule \textbf{Level} & \textbf{Evidence of competency}\\\midrule\endhead
Foundational & Recognizes a weak boundary with prompting; distinguishes a bug from a security impact.\\\midrule
Intermediate & Independently traces representative input to sensitive operations and implements a validated fix.\\\midrule
Advanced & Audits cross-component flow, considers mitigations and exploits prerequisites, and evaluates false positives.\\\midrule
Professional & Produces reproducible, peer-reviewable evidence; prioritizes risk and demonstrates systemic variant prevention.\\\midrule
\bottomrule\end{longtable}\endgroup
\subsection{Finding classification}
\begingroup\footnotesize\setlength\tabcolsep{3pt}
\begin{longtable}{>{\raggedright\arraybackslash}p{\dimexpr0.220\textwidth-6pt\relax}>{\raggedright\arraybackslash}p{\dimexpr0.780\textwidth-6pt\relax}}\toprule
\textbf{Status} & \textbf{Acceptance criteria}\\\midrule\endfirsthead
\toprule \textbf{Status} & \textbf{Acceptance criteria}\\\midrule\endhead
Suspected & Potential weakness with insufficient reachability or security impact evidence; explicitly mark as unconfirmed.\\\midrule
Confirmed & Specific trust violation, reachable attacker conditions and controlled validation or logically sufficient code proof.\\\midrule
False positive & Demonstrated non-reachability or valid defense, supported with executable evidence and control-flow reasoning.\\\midrule
Mitigated & Underlying code issue may remain but documented compensating control blocks the defined abuse path; residual risk required.\\\midrule
Fixed & Root cause is patched, adversarial regression tests pass, and the affected path is analyzed after the change.\\\midrule
\bottomrule\end{longtable}\endgroup
Severity is not the same as exploitability or priority: CVSS characterizes severity; EPSS estimates exploitation likelihood for scored CVEs; impact, exposure, business context, and confidence still require expert evaluation. Do not generate fabricated EPSS values for hypothetical educational defects.

\section{Capstone specification and acceptance criteria}
Assess a deliberately vulnerable Windows service implementing a restricted IPC interface and file operations.

\subsection{Week 15: controlled security assessment}
\begin{itemize}[leftmargin=*,itemsep=2pt]
\item Record signed/authorized scope, target version and prohibited actions.
\item Draw a context-level architecture diagram, system trust boundaries and three explicit security invariants.
\item Identify at least four plausible weaknesses and deeply investigate at least two; zero confirmed findings is acceptable if supported by thorough testing and a justified limitations section.
\item For each finding: source location, attacker control, guard analysis, impact, reproduction steps or logical proof, confidence and precise residual uncertainty.
\item Include at least one independently reviewed false-positive/suspected finding analysis.
\end{itemize}
\subsection{Week 16: remediation and proof}
\begin{itemize}[leftmargin=*,itemsep=2pt]
\item Submit isolated corrective patches and a before/after verification matrix.
\item Implement a regression test on the original trigger and at least two adversarial variants per fixed issue.
\item Run a negative-control experiment to demonstrate that validation is measuring the intended behavior.
\item Produce a SARIF-compatible issue summary or structured triage ledger; do not claim upload to GitHub unless actually performed.
\item Deliver a 10-minute technical defense including remaining risk, systemic prevention and operational handover.
\end{itemize}
Capstone rubric (100 pts): scope and model 15; root-cause and trust evidence 25; controlled validation 20; secure patch 20; regression tests 10; reporting and limitations 10. The week-15 and week-16 submissions are assessed jointly for the course capstone weight (25\%).

\section{Working glossary and reporting schema}
\begingroup\footnotesize\setlength\tabcolsep{3pt}
\begin{longtable}{>{\raggedright\arraybackslash}p{\dimexpr0.240\textwidth-6pt\relax}>{\raggedright\arraybackslash}p{\dimexpr0.760\textwidth-6pt\relax}}\toprule
\textbf{Term} & \textbf{Operational definition}\\\midrule\endfirsthead
\toprule \textbf{Term} & \textbf{Operational definition}\\\midrule\endhead
Source & A location at which data controlled by another trust principal enters an analysis flow.\\\midrule
Sink & A security-sensitive operation, such as code execution, allocation, file access or database query.\\\midrule
Sanitizer / guard & Control that transforms untrusted data or proves a precondition; effectiveness is context-specific.\\\midrule
Reachability & Evidence that conditions connect a relevant input to the vulnerable operation.\\\midrule
TOCTOU & Time-of-check-to-time-of-use hazard from changing a referenced object between validation and action.\\\midrule
False positive & A reported condition shown not to violate a security invariant, with evidence.\\\midrule
Mitigated & Risk blocked by a documented compensating control while the code root cause may persist.\\\midrule
SARIF & Static Analysis Results Interchange Format; structured tool result exchange schema.\\\midrule
Threat model & Explicit assets, actors, capabilities, assumptions, trust boundaries and abuse paths.\\\midrule
\bottomrule\end{longtable}\endgroup
\subsection{Required professional finding fields}
\begin{itemize}[leftmargin=*,itemsep=2pt]
\item Finding ID, affected version/component/file/line, weakness and related classification.
\item Input control, affected trust boundary, source-to-sink or state trace, required preconditions.
\item Observed behavior, controlled validation command/inputs, exact results, negative control.
\item Security impact and confidence, relevant defense-in-depth, limitations and residual risk.
\item Minimal remediation, corrected commit/diff, regression test identifiers, reviewer signoff.
\end{itemize}
\section{Instructor implementation and version control}
\begin{itemize}[leftmargin=*,itemsep=2pt]
\item Provide two instructor-reviewed test fixtures per chapter, one vulnerable and one corrected; lock revision hashes before cohort enrollment.
\item Publish a weekly target specification and evidence submission rubric in the course repository; permit alternative operating systems only when technical outcomes match.
\item Check license entitlement for commercial products and GitHub repository security features; offer local community/free-tool alternatives.
\item Never run uncertain or intentionally unsafe demonstrations on shared services, production accounts or public IP addresses.
\item Repeat version checks each semester. Label contemporary techniques as new/modernized, and distinguish them from historical book concepts.
\item Evaluate learner reasoning through oral code walkthroughs; scanner output without a threat model and relevant code evidence is insufficient.
\end{itemize}
\section{References and verified technical resources}
Status review: 17 October 2026. Stable OWASP ASVS is 5.0.0; the 2025 OWASP Top 10 is released. NIST SP 800-218/SSDF 1.1 is final, while SP 800-218 Rev. 1 / SSDF 1.2 remained listed as draft at review time. Verify moving tool compatibility before deploying laboratories.

\noindent\textbf{Dowd, McDonald, and Schuh. The Art of Software Security Assessment (2007).} Textbook; preface and Chapters 1--18.\par\small\url{https://www.oreilly.com/library/view/the-art-of/0321444426/}\par\normalsize\medskip
\noindent\textbf{MITRE. CWE: Common Weakness Enumeration.} Weakness descriptions and relationships.\par\small\url{https://cwe.mitre.org/}\par\normalsize\medskip
\noindent\textbf{NIST. SP 800-218 (SSDF 1.1, final, 2022).} Secure software development activities.\par\small\url{https://csrc.nist.gov/pubs/sp/800/218/final}\par\normalsize\medskip
\noindent\textbf{NIST. SP 800-218 Rev. 1 (SSDF 1.2, draft as of 2026-10).} Research-only draft; not a replacement for final 1.1.\par\small\url{https://csrc.nist.gov/projects/ssdf/publications}\par\normalsize\medskip
\noindent\textbf{OWASP. ASVS 5.0.0 (stable).} Verifiable application security requirements.\par\small\url{https://owasp.org/www-project-application-security-verification-standard/}\par\normalsize\medskip
\noindent\textbf{OWASP. Top 10:2025 (released).} Web risk awareness categories, not a testing checklist.\par\small\url{https://top10.owasp.org/2025/}\par\normalsize\medskip
\noindent\textbf{OWASP. Web Security Testing Guide (Stable).} Web testing methodologies.\par\small\url{https://wstg.owasp.org/}\par\normalsize\medskip
\noindent\textbf{OWASP. API Security Top 10:2023.} API risk awareness.\par\small\url{https://owasp.org/API-Security/editions/2023/en/0x00-header/}\par\normalsize\medskip
\noindent\textbf{GitHub. CodeQL supported languages.} Language coverage and tool limitations.\par\small\url{https://codeql.github.com/docs/codeql-overview/supported-languages-and-frameworks/}\par\normalsize\medskip
\noindent\textbf{GitHub. CodeQL compiled-language build modes.} CodeQL database production.\par\small\url{https://docs.github.com/en/code-security/concepts/code-scanning/codeql/codeql-for-compiled-languages}\par\normalsize\medskip
\noindent\textbf{FIRST. CVSS 4.0 specification.} Severity characterization, not a probability of exploitation.\par\small\url{https://www.first.org/cvss/v4.0/specification-document}\par\normalsize\medskip
\noindent\textbf{FIRST. Exploit Prediction Scoring System.} Exploit prediction and prioritization.\par\small\url{https://www.first.org/epss/}\par\normalsize\medskip
\noindent\textbf{IETF. RFC 9110 HTTP Semantics.} HTTP behavior shared across versions.\par\small\url{https://www.rfc-editor.org/rfc/rfc9110}\par\normalsize\medskip
\noindent\textbf{IETF. RFC 9114 HTTP/3.} HTTP/3 mapping over QUIC.\par\small\url{https://www.rfc-editor.org/rfc/rfc9114}\par\normalsize\medskip
\noindent\textbf{IETF. RFC 9000 QUIC Transport.} QUIC transport behavior.\par\small\url{https://www.rfc-editor.org/rfc/rfc9000}\par\normalsize\medskip
\noindent\textbf{GitHub. SARIF support for code scanning.} Structured static analysis evidence.\par\small\url{https://docs.github.com/en/code-security/code-scanning/integrating-with-code-scanning/sarif-support-for-code-scanning}\par\normalsize\medskip
\noindent\textbf{Linux. openat2 manual page.} Directory-relative path resolution controls.\par\small\url{https://man7.org/linux/man-pages/man2/openat2.2.html}\par\normalsize\medskip
\noindent\textbf{Microsoft. Windows security documentation.} Security descriptors, access tokens, RPC and ACLs.\par\small\url{https://learn.microsoft.com/en-us/windows/security/}\par\normalsize\medskip
\noindent\textbf{LLVM. AddressSanitizer documentation.} Memory error checking.\par\small\url{https://clang.llvm.org/docs/AddressSanitizer.html}\par\normalsize\medskip
\noindent\textbf{SEI CERT. Secure Coding Standards.} Implementation standards.\par\small\url{https://wiki.sei.cmu.edu/confluence/display/seccode}\par\normalsize\medskip
\end{document}
